\section{Batasan Masalah}

Batasan masalah pada penelitian ini adalah:
\begin{enumerate}
    \item Penelitian dilakukan pada dua ruang-waktu: Schwarzschild dan Schwarzschild 
    Anti-de Sitter

    \item Penelitian melakukan perturbasi medan skalar bermassa dan tanpa massa terhadap 
    kedua ruang-waktu yang dipilih

    \item Penelitian menggunakan metode kombinasi Frobenius dan metode pecahan berlanjut 
    untuk mendapatkan nilai frekuensi eigen quasinormal modes

    \item Penelitian menggunakan parameter momentum sudut $\ell$ = 0, 1, 2; massa medan skalar 
    bernilai satu; dan radius cakrawala sangat kecil dari massa lubang hitam

    \item Penelitian berfokus pada mode fundamental dan mode overtone
\end{enumerate}